\section{Introdução}
\label{introducao}

A análise de imagens de sensoriamento remoto adquiridas antes e depois de um desastre pode contribuir para a caracterização das alterações observáveis na área afetada. Nesse contexto, a comparação temporal permite investigar mudanças em elementos da paisagem, como edificações e cobertura do solo. A disponibilidade de imagens anteriores e posteriores aos eventos e de anotações de danos, exemplificada pelo conjunto xBD, evidencia o interesse em desenvolver métodos computacionais para apoiar esse tipo de análise \citep{gupta2019xbd}. Neste trabalho, esse apoio é entendido como uma tarefa de descrição visual, sem substituir avaliações especializadas ou estabelecer, por si só, a causa das alterações identificadas.

A detecção de mudanças e a descrição textual dessas mudanças são tarefas relacionadas, mas produzem informações distintas. Enquanto a primeira pode indicar as regiões que sofreram alterações, a segunda busca expressar em linguagem natural o que mudou entre duas observações. Assim, além de localizar diferenças, torna-se necessário reconhecer os elementos envolvidos e relacioná-los temporalmente. No sensoriamento remoto, essa formulação é investigada como \textit{remote sensing image change captioning}, que associa pares de imagens a sentenças descritivas das mudanças na cobertura do solo \citep{liu2022}.

Essa tarefa se relaciona à geração automática de legendas para uma única imagem, conhecida como \textit{image captioning}. Nesse problema, a representação visual é utilizada para produzir uma sentença que descreva o conteúdo da cena, aproximando visão computacional e processamento de linguagem natural \citep{vinyals2015}. Em \textit{change captioning}, o foco se desloca do conteúdo isolado para a diferença semântica entre duas imagens. Tal distinção exige separar mudanças relevantes de variações que não correspondem à alteração de interesse, como diferenças de ponto de vista \citep{park2019}.

Ao considerar imagens bitemporais de sensoriamento remoto, essa separação também precisa lidar com fatores como iluminação, sazonalidade e complexidade da cobertura do solo \citep{chang2023}. Uma legenda coerente deve, portanto, representar uma mudança sustentada pelo par de imagens, e não apenas reproduzir uma descrição plausível para uma das cenas. Para o recorte deste trabalho, também será importante verificar se a sentença preserva a ordem entre a observação anterior e a posterior e se evita atribuir ao desastre informações que não estejam visualmente sustentadas.

Entre as abordagens voltadas a esse problema, o Chg2Cap, de \citet{chang2023}, combina extração de características por uma rede convolucional siamesa, processamento atento das representações visuais e geração de legendas baseada em Transformer. A arquitetura Transformer utiliza mecanismos de atenção para modelar relações entre elementos de uma sequência \citep{vaswani2017}. Neste trabalho, o Chg2Cap será adotado como arquitetura de referência, constituindo o ponto de partida para a reprodução experimental e para adaptações incrementais ao domínio de desastres naturais.

\subsection{Problema de pesquisa e justificativa}

A adoção de uma arquitetura existente não assegura que seu desempenho se mantenha em outro domínio. Neste projeto, será necessário investigar tanto a adequação das representações aprendidas quanto a compatibilidade entre as imagens, as anotações disponíveis e as descrições que se pretende gerar. Por isso, a reprodução inicial em um conjunto de referência constitui uma etapa de controle do procedimento experimental, anterior à análise dos efeitos da adaptação.

Outro ponto central é a disponibilidade de legendas de referência. Anotações destinadas a localizar edificações ou classificar danos não equivalem, automaticamente, a sentenças que descrevam mudanças. O xBD, por exemplo, fornece imagens anteriores e posteriores aos eventos, polígonos de edificações e categorias de dano \citep{gupta2019xbd}. Sua menção aqui contextualiza o problema de dados e não representa uma escolha definitiva para os experimentos. A seleção do conjunto de desastres e a estratégia de obtenção ou adequação das legendas ainda integram o planejamento metodológico.

Diante disso, formula-se a seguinte questão de pesquisa: \textit{em que medida uma abordagem baseada no Chg2Cap consegue descrever mudanças observáveis em imagens bitemporais de sensoriamento remoto de cenários de desastres naturais, e quais limitações se manifestam nessa adaptação?} A investigação se justifica pela possibilidade de avaliar a descrição textual como uma forma complementar de apresentar as alterações visuais, examinando seus limites sem pressupor utilidade operacional ou superioridade em relação a outras abordagens.

\subsection{Objetivos}

O objetivo geral é adaptar e avaliar uma abordagem baseada no Chg2Cap para gerar descrições textuais das mudanças observadas em imagens bitemporais de sensoriamento remoto em cenários de desastres naturais.

Para alcançar esse objetivo, estabelecem-se os seguintes objetivos específicos:

\begin{enumerate}
    \item Reproduzir e documentar o funcionamento do Chg2Cap no conjunto de referência LEVIR-CC, registrando as configurações, os procedimentos de avaliação e as legendas geradas.
    \item Selecionar e preparar dados de cenários de desastres naturais, considerando o pareamento temporal das imagens e a disponibilidade ou viabilidade de obtenção de legendas de referência.
    \item Adequar o processamento dos dados e o procedimento experimental ao domínio selecionado, preservando inicialmente a arquitetura de referência.
    \item Executar experimentos de treinamento e avaliação, com registro das configurações e separação dos dados utilizados nas diferentes etapas.
    \item Analisar os resultados por meio de métricas de geração de legendas e de inspeção qualitativa, identificando erros de conteúdo, mudanças inexistentes, inversões temporais e omissões relevantes.
\end{enumerate}

\subsection{Delimitação e contribuição esperada}

O escopo compreende pares de imagens da mesma área em dois momentos e a geração de uma descrição das mudanças observáveis entre eles. O LEVIR-CC, apresentado por \citet{liu2022}, será utilizado como referência inicial, sem que seus resultados sejam interpretados como evidência de desempenho em desastres naturais. A contribuição esperada é de natureza experimental e aplicada: documentar a reprodução, examinar a adaptação de dados e analisar capacidades e limitações do Chg2Cap no domínio selecionado. Não se pressupõe a criação de uma nova arquitetura, nem se antecipam resultados ainda não obtidos.

\subsection{Organização do trabalho}

O texto está organizado em cinco seções. Após esta introdução, a Seção~\ref{referencial} reúne os fundamentos teóricos. A Seção~\ref{desenvolvimento} será dedicada aos dados e aos procedimentos experimentais. A Seção~\ref{resultados} apresentará os resultados e sua análise, e a Seção~\ref{conclusoes} reunirá as conclusões, limitações e possibilidades de continuidade.
